\section*{第 \ref{ch:g12:matrix} 章 --- 矩阵与图}

\begin{solution}{exo:g12:matrix:1}
\[
A + B = \begin{pmatrix} 3 & 2\\ 1 & 2\end{pmatrix}, \quad
AB = \begin{pmatrix} 4 & 2\\ 1 & 1\end{pmatrix}, \quad
BA = \begin{pmatrix} 2 & 4\\ 1 & 3\end{pmatrix}, \quad
A^2 = \begin{pmatrix} 1 & 4\\ 0 & 1\end{pmatrix}.
\]
注 $AB \neq BA$。
\end{solution}

\begin{solution}{exo:g12:matrix:2}
$\det A = 6 - 5 = 1 \neq 0$:
$A^{-1} = \begin{pmatrix} 3 & -5\\ -1 & 2 \end{pmatrix}$.
$\det B = 12 - 12 = 0$：$B$ 不可逆。
\end{solution}

\begin{solution}{exo:g12:matrix:3}
方程组为 $AX = Y$，其中 $A$ 与\cref{exo:g12:matrix:2}相同，且
$Y = \begin{pmatrix} 1\\ 2\end{pmatrix}$：
\[
X = A^{-1}Y = \begin{pmatrix} 3 & -5\\ -1 & 2\end{pmatrix}
\begin{pmatrix} 1\\ 2\end{pmatrix}
= \begin{pmatrix} -7\\ 3\end{pmatrix}:
\qquad x = -7,\ y = 3 .
\]
\end{solution}

\begin{solution}{exo:g12:matrix:4}
\emph{1.} $N^2 = \begin{pmatrix} 0&1\\0&0\end{pmatrix}
\begin{pmatrix} 0&1\\0&0\end{pmatrix} = 0$，且 $I$ 与每个
矩阵都可交换。

\emph{2.} 由于两项可交换，可用二项式定理，且
所有含 $N^2$ 的项均为零：
\[
A^k = (2I + N)^k = 2^k I + k\,2^{k-1} N
= \begin{pmatrix} 2^k & k\,2^{k-1}\\ 0 & 2^k \end{pmatrix}.
\]
当 $k = 2$ 时检验：
$A^2 = \begin{pmatrix} 2&1\\0&2\end{pmatrix}^2
= \begin{pmatrix} 4&4\\0&4\end{pmatrix}$，而公式给出
$2^2 = 4$, $2 \times 2 = 4$. \checkmark
\end{solution}

\begin{solution}{exo:g12:matrix:5}
\emph{归纳法。}当 $n = 1$ 时：
$A^1 = \begin{pmatrix} 0&1\\1&1\end{pmatrix}
= \begin{pmatrix} F_0 & F_1\\ F_1 & F_2\end{pmatrix}$。设公式对
$n$ 成立，则
\[
A^{n+1} = A^n A
= \begin{pmatrix} F_{n-1} & F_n\\ F_n & F_{n+1}\end{pmatrix}
\begin{pmatrix} 0 & 1\\ 1 & 1\end{pmatrix}
= \begin{pmatrix} F_n & F_{n-1} + F_n\\ F_{n+1} & F_n + F_{n+1}\end{pmatrix}
= \begin{pmatrix} F_n & F_{n+1}\\ F_{n+1} & F_{n+2}\end{pmatrix}.
\]
\emph{恒等式。}对 $2\times2$ 矩阵，展开可知
$\det(MN) = \det M \det N$；由此
可得 $\det(A^n) = (\det A)^n = (-1)^n$，且
$\det A^n = F_{n-1}F_{n+1} - F_n^2$。（这就是\emph{卡西尼恒等式}。）
\end{solution}

\begin{solution}{exo:g12:matrix:6}
\emph{1.} 顶点按 $1, 2, 3$ 排序：
\[
M = \begin{pmatrix} 0&1&1\\ 0&0&1\\ 1&0&0\end{pmatrix}, \quad
M^2 = \begin{pmatrix} 1&0&1\\ 1&0&0\\ 0&1&1\end{pmatrix}, \quad
M^3 = \begin{pmatrix} 1&1&1\\ 1&0&1\\ 1&0&1 \end{pmatrix}.
\]

\emph{2.} $\bigl(M^3\bigr)_{11} = 1$：长为 $3$ 的闭途径中，
从顶点 $1$ 出发的恰有一条，即 $1 \to 2 \to 3 \to 1$。（途径
$1 \to 3 \to 1$ 的长度仅为 $2$，而 $1 \to 3$ 后接 $3\to1$ 再接
$1\to3$ 则终止于 $3$。）
\end{solution}

\begin{solution}{exo:g12:matrix:7}
\emph{1.} $a_{n+1} = 0.8a_n + 0.3b_n$, $b_{n+1} = 0.2a_n + 0.7b_n$:
$M = \begin{pmatrix} 0.8 & 0.3\\ 0.2 & 0.7\end{pmatrix}$.

\emph{2.} $MX = X$ 给出 $0.8a + 0.3b = a$，\emph{即} $0.3b = 0.2a$，故
$b = \frac23 a$；结合 $a + b = 1$：$a = 0.6$，$b = 0.4$。

\emph{3.} 利用 $b_n = 1 - a_n$：
$a_{n+1} = 0.8a_n + 0.3(1 - a_n) = 0.5a_n + 0.3$，故
\[
c_{n+1} = a_{n+1} - 0.6 = 0.5a_n + 0.3 - 0.6 = 0.5(a_n - 0.6) = 0.5\,c_n .
\]
故 $c_n = 0.5^n c_0 \to 0$：$a_n \to 0.6$ 且 $b_n \to 0.4$，无论
初始分布如何。
\end{solution}

\begin{solution}{exo:g12:matrix:8}
\emph{1.} $\det P = -2$，故
$P^{-1} = -\frac12\begin{pmatrix} -1 & -1\\ -1 & 1\end{pmatrix}
= \frac12\begin{pmatrix} 1 & 1\\ 1 & -1\end{pmatrix}$。则
\[
AP = \begin{pmatrix} 4 & 2\\ 4 & -2 \end{pmatrix}, \qquad
D = P^{-1}AP = \frac12\begin{pmatrix} 1&1\\1&-1\end{pmatrix}
\begin{pmatrix} 4&2\\4&-2\end{pmatrix}
= \begin{pmatrix} 4 & 0\\ 0 & 2\end{pmatrix}.
\]

\emph{2.} 由 $A = PDP^{-1}$，经归纳可得
$A^n = PD^nP^{-1}$，其中
$D^n = \begin{pmatrix} 4^n & 0\\ 0 & 2^n\end{pmatrix}$，故
\[
A^n = P D^n P^{-1}
= \begin{pmatrix} 4^n & 2^n\\ 4^n & -2^n\end{pmatrix}\cdot
\frac12\begin{pmatrix} 1&1\\1&-1\end{pmatrix}
= \frac12\begin{pmatrix} 4^n + 2^n & 4^n - 2^n\\
4^n - 2^n & 4^n + 2^n\end{pmatrix}.
\]
将 $A^n$ 作用于 $X_0 = \begin{pmatrix} u_0\\v_0\end{pmatrix}$，
恰好重现\cref{ex:g12:matrix:coupled}中的公式。
\end{solution}

\begin{solution}{pb:g12:matrix:1}
\textbf{1.} $AB = \begin{pmatrix} 2 & 1\\ 4 & 3\end{pmatrix}$
而 $BA = \begin{pmatrix} 3 & 4\\ 1 & 2\end{pmatrix}$：矩阵
乘法不满足交换律 --- $B$ 右乘时交换列，
左乘时交换行。

\textbf{2.} 行列式 $6 - 5 = 1$：逆矩阵
为 $\begin{pmatrix} 3 & -1\\ -5 & 2\end{pmatrix}$。将其作用于
$\binom{4}{7}$：$x = 12 - 7 = 5$，$y = -20 + 14 = -6$。

\textbf{3.} $N^2 = 0$。于是由归纳法得 $(I + N)^n = I + nN$：
$(I + nN)(I + N) = I + (n+1)N + nN^2 = I + (n+1)N$.

\textbf{4.} $A = \begin{pmatrix} 0&1&1\\ 1&0&1\\ 1&1&0
\end{pmatrix}$，且 $A^3$ 的对角元均为 $2$：从每个
顶点出发，恰有两条长为 $3$ 的闭途径（三角形
按顺时针或逆时针方向绕行） --- 计数定理
的应用。

\textbf{5.} $D^n = \begin{pmatrix} 2^n & 0\\ 0 & 2^{-n}
\end{pmatrix}$：一个方向迅速增大，另一个方向衰减至零 ---
对角线上的两种“命运”是相互独立的等比数列。

\textbf{6.} $F^2 = \begin{pmatrix} 2 & 1\\ 1 & 1\end{pmatrix}$,
$F^3 = \begin{pmatrix} 3 & 2\\ 2 & 1\end{pmatrix}$,
$F^4 = \begin{pmatrix} 5 & 3\\ 3 & 2\end{pmatrix}$：斐波那契
数无处不在；猜想如题所述。

\textbf{7.} 若 $F^n = \begin{pmatrix} F_{n+1} & F_n\\ F_n &
F_{n-1}\end{pmatrix}$，则
\[
F^{n+1} = F^n F =
\begin{pmatrix} F_{n+1} + F_n & F_{n+1}\\
F_n + F_{n-1} & F_n \end{pmatrix}
= \begin{pmatrix} F_{n+2} & F_{n+1}\\ F_{n+1} & F_n
\end{pmatrix} :
\]
归纳步；基础情形 $n = 1$ 即 $F$ 本身，这里用到的
约定 $F_0 = 0$（它将递推关系向前延拓）。

\textbf{8.} $\det F = -1$，故
$\det(F^n) = (\det F)^n = (-1)^n$；而直接计算得
$\det(F^n) = F_{n+1}F_{n-1} - F_n^2$：卡西尼恒等式，一行即可。（
$2 \times 2$ 行列式的乘积法则是一段愉快的
五分钟展开计算。）

\textbf{9.} $F^m F^n$ 的右上角元素：
$F_{m+1}F_n + F_m F_{n-1}$；$F^{m+n}$ 的右上角元素为
$F_{m+n}$。取 $m = n = 3$：
$F_4 F_3 + F_3 F_2 = 3 \times 2 + 2 \times 1 = 8 = F_6$.

\textbf{10.} 当 $k = 1$ 时显然成立。若 $F_n \mid F_{kn}$，
取 $m = kn$ 的加法公式：
$F_{(k+1)n} = F_{kn+1}F_n + F_{kn}F_{n-1}$：两项均为
$F_n$ 的倍数。故 $F_n \mid F_{kn}$ 对所有 $k$ 成立。验证：
$F_3 = 2$ 整除 $F_6 = 8$ 及 $F_9 = 34$。

\textbf{11.} $F^{100} = F^{64} F^{32} F^4$：七次平方运算
（$F^2, F^4, \dots, F^{64}$）外加两次组合运算 --- 共
九次乘法，而非九十九次。这正是
古埃及抄写员的“倍乘表”技巧，由数推广
到矩阵：将 $100$ 写成二进制，
然后乘上所需的各次倍幂。

\textbf{12.} $0.8 + 0.2 = 1$ 与 $0.4 + 0.6 = 1$：明天
必定是\emph{某种}天气 --- 每一列都是一个完整的
概率分布，因此概率守恒。

\textbf{13.} 明天：$\binom{0.8}{0.2}$。后天：
$M\binom{0.8}{0.2} = \binom{0.72}{0.28}$.

\textbf{14.} $Mv = v$，其中 $v = \binom{s}{r}$，$s + r = 1$：
$0.8s + 0.4r = s$ 得 $0.4r = 0.2s$，$s = 2r$：
$v = \binom{2/3}{1/3}$。长期来看，三天中有两天
是晴天 --- 无论今天天气如何。

\textbf{15.} 自 $\binom01$ 出发，晴天分量为 $0.4$，$0.56$，
$0.624$，$0.6496$；与 $\frac23$ 的差距分别为 $0.267$，$0.107$，
$0.043$，$0.017$ --- 每一步都将差距恰好乘以
$0.4$（即该机器的第二特征值）：等比收敛
至稳态。

\textbf{16.} 列（依次取自 $A$、$B$、$C$）：
$P = \begin{pmatrix} 0 & 0 & 1\\ \frac12 & 0 & 0\\
\frac12 & 1 & 0\end{pmatrix}$。稳态：$v_A = v_C$，
$v_B = \frac{v_A}{2}$，$v_C = \frac{v_A}{2} + v_B$；总和为
$1$：$v = \left(\frac25, \frac15, \frac25\right)$。排名：
$A$ 与 $C$ 并列第一，$B$ 居末。

\textbf{17.} $C$ 接收\emph{全部}来自 $B$ 的流量及
$A$ 流量的一半，并将所有流量导回 $A$：稳态
度量的是冲浪者\emph{停留时间}之所在，而非
有多少链接指向它 --- 一个来自热门网页的链接
胜过数个来自冷门网页的链接。这种递归式加权
正是谷歌的奠基思想；阻尼用于处理蜘蛛陷阱
与死胡同。

\textbf{18.} 由 $A^n = P D^n P^{-1}$ 可得
$u_n = \frac{4^n + 2^n}{2}$（及
$v_n = \frac{4^n - 2^n}{2}$）。验证：$u_1 = 3$，$u_2 = 10$，
$u_3 = 36$；直接计算：$(1,0) \to (3,1) \to (10,6) \to
(36, 28)$：吻合。

\textbf{19.} 对角化即换到一组坐标，在其中
耦合系统分解为相互独立的等比
数列 --- 每个特征值各跑各的赛程。马尔可夫
链的稳态即特征值 $1$ 的特征向量，而
第 15 题中的收敛速率即下一个特征值：
天气机器始终是一个特征值的故事。

\textbf{20.} 记账：一个系统即一个矩阵方程，
由一个逆矩阵求解。组合学：邻接矩阵的幂
计数途径、链接与连接。演化：机器的幂
将各状态带向各自的归宿，而
特征方向（斐波那契的黄金方向、天气的
稳态、网页的排序向量）即为归宿。
线性代数在大学分册中，正是研究
这类问题的学科。
\end{solution}
